\chapter{Toàn bộ cỗ máy}
% full version: chapters/10_synthesis.tex

\pencilfig{10}{\pencilart{10}}{Toàn bộ cỗ máy: bus dữ liệu, bộ phanh bên trong nó và bốn kênh radio điều khiển ở phía trên.}

Chúng ta đã lần lượt xét từng loại nơ-ron và hỏi nó thêm gì cho cỗ máy. Đã đến lúc lắp các linh kiện lại với nhau.

Kết quả là một cỗ máy ba tầng. Ở dưới là mạng điện thoại khổng lồ của các nơ-ron \nt{GLUT}: dữ liệu chạy trên đó. Bên cạnh là các nơ-ron ức chế \nt{GABA}: chúng điều khiển dòng chảy, quyết định cái gì được đi qua, khi nào và to đến đâu. Và phía trên tất cả là vài đài phát thanh, mỗi đài có một núm riêng. \nt{DOPA} cho biết nên giữ lại những thay đổi liên kết nào. \nt{SERO} giữ mức chung và, theo giả thuyết, tầm xa của sự kiên nhẫn. \nt{NORA} chọn giữa tìm kiếm và quyết định. \nt{ACET} --- giữa ghi và đọc. Thực ra mỗi đài phát thanh là một bó kênh nhỏ chứ không phải một dây dẫn, nhưng dù vậy số kênh vẫn chỉ là một nhúm cho tám mươi sáu tỷ nơ-ron.

\section*{Một câu chuyện}

Hãy theo dõi một sự kiện đi qua toàn bộ cỗ máy. Con vật từ lâu đã biết: nhấn cần A --- sẽ có nước ép. Một ngày thế giới thay đổi: A không còn tác dụng, còn nước ép giờ đến từ cần B, cái cần mà con vật hầu như chưa dùng.

Nhấn A mà nước ép không tới --- dopamine đáp lại bằng một khoảng lặng, và các liên kết đã chọn A yếu đi. Các khoảng lặng lặp lại, sai số trở nên lớn hơn thường lệ --- noradrenaline, theo giả thuyết, báo rằng thế giới đã đổi. Núm độ tương phản hạ xuống, lựa chọn trở nên mềm hơn, và nhiễu thường xuyên hơn đẩy con vật thử B. Acetylcholine chuyển mạng sang chế độ ghi. Lần thử B mang lại nước ép bất ngờ --- dopamine bùng lên, các liên kết đã chọn B mạnh lên. Sau vài lần thử, kỳ vọng bắt kịp thế giới mới, bất ngờ hết, các núm trở về vị trí cũ.

Hãy để ý: không núm nào biết các núm khác đang làm gì. Mỗi núm phản ứng với dấu hiệu của riêng nó, và trình tự hành động tự hình thành, như trong một mạch bất đồng bộ, nơi mỗi khối kích hoạt khi các đầu vào của nó đã sẵn sàng. Việc từng núm riêng lẻ hành xử như vậy phần lớn đã được đo. Còn việc chúng phối hợp nhịp nhàng với nhau thì hiện vẫn là giả thuyết.

\section*{Nó khác chiếc máy tính xách tay của bạn ở đâu}

Máy tính xách tay có một nhịp đồng hồ chung --- bộ não thì không. Các phần tử của máy tính xách tay đáng tin cậy --- còn khớp thần kinh làm mất hơn một nửa số tín hiệu. Bộ nhớ của máy tính xách tay tách khỏi bộ xử lý --- trong bộ não nó nằm ngay trong các liên kết. Máy tính xách tay học (nếu có học) bằng lan truyền ngược sai số --- bộ não học bằng các quy tắc cục bộ cộng với một tín hiệu quảng bá ``tốt hơn mong đợi''. Và điều chính: máy tính xách tay tốn hàng chục oát cho riêng bộ xử lý, còn bộ não --- hai mươi oát cho tất cả.

\section*{Điều chúng ta chưa biết}

Chúng ta không biết vỏ não dùng thời điểm chính xác của tiếng tách tới mức nào. Không biết nó điều chỉnh hàng tỷ liên kết trong các mạch nhiều tầng ra sao, khi thầy giáo chỉ có một cho tất cả. Chưa hiểu hết các đài phát thanh truyền đi những gì. Không biết các phần xa nhau của bộ não phối hợp công việc thế nào khi không có nhịp chung. Và đến nay chưa ai biết cách chế tạo một cỗ máy làm được điều tương tự với hai mươi oát.

Phần tiếp theo của cuốn sách vượt ra ngoài cái lõi này: tới các đầu vào và đầu ra của cỗ máy, tới việc gỡ lỗi nó, tới giấc ngủ và tới những cỗ máy làm từ nơ-ron sống trên chip.

\fullversion{10}
